\section*{第 \ref{ch:b3:pericyclic} 章——周环反应}

\begin{solution}{exo:b3:pericyclic:1}
第尔斯–阿尔德：[4+2] 环加成。环丁烯开环：电环化。科普重排：[3,3]
σ 迁移。砜烯失去 $\ce{SO2}$：螯键反应。臭氧与烯烃：1,3-偶极
环加成（[3+2]）。$[1,5]$-H 迁移：σ 迁移。
\end{solution}

\begin{solution}{exo:b3:pericyclic:2}
己三烯（6 个电子）：热反应为对旋，光化学反应为顺旋。
丁二烯（4）：热反应为顺旋，光化学反应为对旋。
\end{solution}

\begin{solution}{exo:b3:pericyclic:3}
4 个电子，热反应：顺旋开环。两个甲基位于环的相反两面，二者都向外转动：
（\textit{E},\textit{E}）-己-2,4-二烯。
\end{solution}

\begin{solution}{exo:b3:pericyclic:4}
同面 [2+2] 反应含 $4q$ 个电子：基态时，一个乙烯的 HOMO 与另一个乙烯的 LUMO
重叠，其两端分别为成键与反键。
激发态时，一个分子中单占据的 $\pi^*$ 轨道
充当 HOMO，两端均可匹配。
\end{solution}

\begin{solution}{exo:b3:pericyclic:5}
保留镜面。己三烯：$\psi_1$ 为 S，$\psi_2$ 为 A，$\psi_3$ 为 S，均已占据；$\psi_4$ 为 A，$\psi_5$ 为 S，$\psi_6$ 为 A，均为空。
环己二烯，依次为：$\sigma$ 为 S，$\pi_1$ 为 S，$\pi_2$ 为 A，均已占据；$\pi_3^*$ 为 S，$\pi_4^*$ 为 A，$\sigma^*$ 为 A，均为空。按对称性
顺序对应，$\psi_1 \to \sigma$，$\psi_2 \to \pi_2$，$\psi_3 \to \pi_1$：占据对占据，对旋闭环
为热允许。
\end{solution}

\begin{solution}{exo:b3:pericyclic:6}
6 个电子，光照下为顺旋，因此两个末端甲基位于
相反两面：\textit{trans}-5,6-二甲基环己-1,3-二烯（热反应得到
\textit{cis} 异构体）。
\end{solution}

\begin{solution}{exo:b3:pericyclic:7}
C5 上的氢迁移到二烯的 C1 位，该位在环上与 C5 相邻，
经六电子过渡态（$\sigma2_s + \pi4_s$）同面进行：一次迁移
得到 1-甲基环戊二烯，再一次迁移得到 2-甲基环戊二烯。五碳链使
氢留在同一面，小环本就如此约束。
\end{solution}

\begin{solution}{exo:b3:pericyclic:8}
$\pi2_s$（烯烃）$+ \pi2_a$（烯酮 C=C，在相反两面受进攻，经由
与之垂直的 C=O $\pi^*$）：$(4q + 2)_s$ 组分计数为 1（即烯烃），为奇数：
热允许。
\end{solution}

\begin{solution}{exo:b3:pericyclic:9}
在椅式构象中，乙烯基的 C=C（原子 1--2）、氧（3）、$\ce{OCH2}$ 中的
碳（4）和烯丙基的 C=C（5--6）；O--C 键断裂，C1--C6 键形成：
产物为戊-4-烯醛，
$\ce{OHC-CH2-CH2-CH=CH2}$。
\end{solution}

\begin{solution}{exo:b3:pericyclic:10}
8 个电子属 $4q$ 类型；含一个异面组分时，为默比乌斯拓扑，
具有 $4q$ 个电子，呈芳香性：过渡态为芳香性，并使
反应热允许（与辛四烯的顺旋闭环相同）。
\end{solution}

\begin{solution}{exo:b3:pericyclic:11}
系数最大的原子配对，使 N3 与末端碳成键：
1,4-二取代三唑。无催化剂时，系数相差不大，
1,4 和 1,5 两种异构体均生成；铜（I）使反应分步进行，经由
炔化铜中间体，且仅得 1,4 异构体。
\end{solution}

\begin{solution}{exo:b3:pericyclic:12}
氢桥接 $j$ 原子自由基 SOMO 的两端，$k = (j + 1)/2$。其末端
系数分别正比于 $\sin(k\pi/(j + 1))$ 与 $(-1)^{k+1}\sin(k\pi/(j + 1))$。当 $j = 7$、$k = 4$ 时，符号相反，故
两端的同号瓣分处相反两面：异面。当
$j = 5$、$k = 3$ 时，同号瓣在同一面：同面。
\end{solution}

\begin{solution}{pb:b3:pericyclic:1}
\textbf{1.} B 环的 C9--C10 $\sigma$ 键；与 5,7-二烯一起，构成己三烯，
C10=C5--C6=C7--C8=C9。

\textbf{2.} 6 个电子（两个 $\pi$ 键与一个 $\sigma$ 键）：电环化开环。

\textbf{3.} 热条件下，环己二烯与其开环产生的己三烯之间只能
缓慢达到平衡，需越过较高的势垒，且环状一侧更稳定；
二烯吸收紫外光，其激发态
在皮秒内即开环。

\textbf{4.} 顺旋：激发态中，三烯体系的单占据 LUMO
（$k = 4$）两端的符号相反。

\textbf{5.} C6=C7 键来自环：其两侧链的延续部分 C5 与 C8，
位于其同一侧，并保持不变：\textit{Z}。

\textbf{6.} 热开环只允许对旋，其势垒远超
体温所能提供的能量，且会沿能量上坡方向进行，生成较不稳定的三烯。

\textbf{7.} 由 C19（C10 上的甲基）到 C9：$[1,7]$-H σ 迁移。

\textbf{8.} 8 个：三烯的 6 个 $\pi$ 电子与 C--H 键的 2 个电子。

\textbf{9.} $\sigma2 + \pi6$。若两者均为同面（s），则各有一个 $(4q + 2)_s$ 组分，
总数为偶，禁阻。若三烯取异面，即 $\sigma2_s + \pi6_a$，计数组分为一个，数为奇，
允许。

\textbf{10.} 7 个原子构成螺旋状、\textit{Z}构型的链，使氢能从一面
转到另一面；三原子链无法扭转到这样的程度。

\textbf{11.} 默比乌斯型（含一个异面组分），8 个电子，即 $4q$：芳香性的。

\textbf{12.} 该反应热允许；其势垒在体温下可被缓慢跨越。

\textbf{13.} 光照下 6 个电子：又为顺旋，但末端可以朝另一方向旋转，
使 H9 与 C19 甲基位于另一面，生成闭环
的立体异构体，即光甾醇。

\textbf{14.} 中央键为 \textit{E} 构型时，C19 与 C9 位于链的两侧，
相距甚远，氢原子无法到达。

\textbf{15.} 前维生素 D$_3$ 自身也吸收光能，并转化为光甾醇和速甾醇，
二者均不能生成维生素 D；最终达到光稳态混合物，使
生成量受到限制。

\textbf{16.} 开环发生在激发态，耗时数飞秒至数皮秒；
这一迁移为热反应，势垒为 \qty{85}{kJ/mol}。

\textbf{17.} $t_{1/2} = \ln 2/k = 0.693/\qty{0.023}{h^{-1}} = \qty{30}{h}$.

\textbf{18.} $1 - \eu^{-0.023 \times 24} = 0.42$.

\textbf{19.} $\ln 10/k = 2.303/\qty{0.023}{h^{-1}} = \qty{100}{h}$.

\textbf{20.} $k(\qty{293.15}{K}) = k(\qty{310.15}{K})\,\eu^{-(E_a/R)(1/293.15 - 1/310.15)} = 0.023 \times \eu^{-1.91} = \qty{3.4e-3}{h^{-1}}$.

\textbf{21.} $2.303/\qty{3.4e-3}{h^{-1}} = \qty{680}{h}$，约 28 天。

\textbf{22.} 热步骤在 \qty{37}{\celsius} 时比在 \qty{20}{\celsius} 时快 7 倍，因此
晨间日照下生成的前维生素，几天之内即可完成转化。

\textbf{23.} 甾体骨架带有众多手性碳原子，
其结构已由甾醇提供；随后光照与温热，
与在皮肤中一样，依次完成两步周环反应。

\textbf{24.} 在体温下，90\,\% 的前维生素 D$_3$ 约在
\qty{100}{h}（四天）内转化为维生素 D$_3$。
\end{solution}
